%==============================================================
%  Chapter 2 -- Scalable K-Step Markov-Chain Recommendation
%                via Matrix Factorization
%  Source: Semester-1 mid-term deck (Matrix Factorization).
%==============================================================
\chapter{Scalable \texorpdfstring{$K$}{K}-Step Markov-Chain Recommendation via Matrix Factorization}
\label{ch:markov-factorization}

This chapter develops the first line of work undertaken for the IndiaMART
recommendation problem: modelling product-to-product navigation as a Markov
chain and using its $K$-step transition behaviour to generate
recommendations. The central difficulty is one of scale. The transition
matrix is of the order of tens of millions of products on a side, and the
naive computation of its $K$-step powers is infeasible both in time and in
memory. We motivate a matrix-factorization reformulation that expresses the
$K$-step transition matrix in terms of two thin factors, analyse the memory
it saves, and then examine three candidate factorization schemes ---
Alternating Least Squares (ALS), Non-negative Matrix Factorization (NMF), and
Weighted NMF --- comparing them in the light of the non-negativity and
sparsity structure that the problem imposes.

%--------------------------------------------------------------
\section{Problem Statement}
\label{sec:mf-problem}

The starting point is a product-to-product association matrix
$M \in \mathbb{R}^{N \times N}$, where each entry $M_{ij}$ encodes the
strength of the transition from product $i$ to product $j$ --- in our setting,
the (empirical) probability that a user who interacts with product $i$ next
interacts with product $j$. The goal is to build a \emph{$K$-step Markov-chain
recommender}: rather than recommending only the immediate successors of a
product, we follow the chain forward $K$ steps so that recommendations capture
multi-hop navigation patterns across the catalogue.

Two structural properties of $M$ dominate the design:

\begin{itemize}
  \item \textbf{Scale.} The matrix is approximately
        $2.3~\text{crore} \times 2.3~\text{crore}$, i.e.\ roughly
        $N \approx 2.3 \times 10^{7}$ on each axis.
  \item \textbf{Sparsity.} The density is approximately $10^{-4}$; the
        overwhelming majority of product pairs never co-occur in a session.
\end{itemize}

The combination is what makes the problem hard. While $M$ itself is sparse and
storable, its powers are not: as the chain is advanced through successive
steps, paths between products proliferate and the density of $M^{k}$ grows
\emph{exponentially} with $k$. Within a few steps the matrix becomes
effectively dense, and materialising it would require an amount of memory that
no commodity machine can provide --- in practice leading to buffer overflow
and out-of-memory failure. The remainder of the chapter is concerned with
computing the $K$-step transition behaviour \emph{without} ever forming a
dense $M^{k}$.

%--------------------------------------------------------------
\section{Background}
\label{sec:mf-background}

\subsection{Recommender Systems}
A recommender system is a data-driven system that surfaces items relevant to a
user by exploiting patterns and relationships latent in interaction data.
Broadly, three families are distinguished. \emph{Content-based} methods
recommend items similar to those a user has already engaged with, using item
attributes. \emph{Collaborative filtering} methods rely instead on the
behaviour of the user population, recommending items favoured by users with
similar histories. \emph{Hybrid} approaches combine the two to offset their
individual weaknesses. Recommender systems are now ubiquitous in e-commerce,
streaming, and social platforms. They are nonetheless beset by recurring
challenges: the \emph{cold-start} problem for new users or items, \emph{data
sparsity} when interactions are few relative to the catalogue, and
\emph{scalability} as catalogues grow to the scale considered here. The work
in this thesis sits squarely against the last two of these.

\subsection{Markov Chains and \texorpdfstring{$K$}{K}-Step Transitions}
\label{sec:markov}
A Markov chain is a stochastic model of a sequence of events in which the
probability of the next event depends only on the present state and not on the
history that preceded it (the \emph{memoryless}, or first-order Markov,
property). The dynamics are captured by a transition matrix $M$, whose entry
$M_{ij}$ gives the probability of moving from state $i$ to state $j$ in a
single step.

The key fact we exploit is how multi-step transitions compose. The
probability of moving from state $i$ to state $j$ in exactly $k$ steps is the
$(i,j)$ entry of the $k$-th matrix power:
\begin{equation}
  \label{eq:kstep}
  \bigl[M^{k}\bigr]_{ij}
    \;=\; \Pr\!\left(\text{state } j \text{ after } k \text{ steps}
                     \mid \text{state } i \text{ now}\right).
\end{equation}
Equation~\eqref{eq:kstep} is exactly the object a $K$-step recommender needs:
ranking the row $\bigl[M^{K}\bigr]_{i,\,\cdot}$ yields the products most likely
to be reached from product $i$ within $K$ steps. The computational obstacle of
Section~\ref{sec:mf-problem} is now precise --- we must read off rows of
$M^{K}$ without forming $M^{K}$ explicitly.

%--------------------------------------------------------------
\section{Proposed Solution: Factorized \texorpdfstring{$K$}{K}-Step Powers}
\label{sec:mf-solution}

Direct iteration of a $2.3\times10^{7}$ square matrix to the $K$-th power is
impossible at this scale. Our proposal is to replace $M$ by a low-rank
factorization and to push the matrix power onto the (tiny) factors.

Suppose $M$ admits a rank-$r$ factorization
\begin{equation}
  \label{eq:factorize}
  M \;\approx\; U V^{\top},
  \qquad U \in \mathbb{R}^{N \times r},\;\; V \in \mathbb{R}^{N \times r},
\end{equation}
with $r \ll N$ (in our experiments $r$ is of the order of $100$). The two
factors $U$ and $V$ are thin: their storage is linear in $N$ rather than
quadratic. The power of this representation is that the $K$-step matrix can be
expressed entirely through a small $r \times r$ core. Substituting
\eqref{eq:factorize} and using associativity,
\begin{align}
  M^{k}
    &\approx \underbrace{(U V^{\top})(U V^{\top}) \cdots (U V^{\top})}_{k\ \text{factors}}
     \nonumber\\[2pt]
    &= U \,\bigl(V^{\top} U\bigr)^{k-1}\, V^{\top}
     \;=\; U\, W^{\,k-1}\, V^{\top},
  \label{eq:factored-power}
\end{align}
where we have defined the \emph{core matrix}
\begin{equation}
  \label{eq:core}
  W \;=\; V^{\top} U \;\in\; \mathbb{R}^{r \times r}.
\end{equation}
Equation~\eqref{eq:factored-power} is the heart of the approach: the only
power that must be computed is $W^{\,k-1}$, an $r\times r$ matrix with
$r \approx 100$, which is trivial. The large factors $U$ and $V$ are touched
only by the final, lazy left- and right-multiplications, and a single
recommendation row $\bigl[M^{K}\bigr]_{i,\cdot}
= U_{i,\cdot}\,W^{\,K-1}\,V^{\top}$ can be evaluated on demand without
materialising the full $N \times N$ result.

\begin{remark}
The mid-term presentation wrote the relation compactly as
$M^{k} = U\,(W^{k})\,V^{\top}$. The exact algebra in
\eqref{eq:factored-power} gives the exponent $k-1$ rather than $k$, since one
factor of $V^{\top}U$ is consumed by the outer $U$ and $V^{\top}$. The
distinction does not affect the memory argument below.
\end{remark}

\subsection{Memory Analysis}
\label{sec:mf-memory}
The saving is best appreciated with the concrete figures of the IndiaMART
setting. Take $N = 2.3 \times 10^{7}$, rank $r = 100$, step count $K = 6$, and
double-precision storage at $8$ bytes per entry.

If $M^{k}$ were materialised densely --- which, as noted, it effectively
becomes after a few steps --- it would occupy
\begin{equation}
  N^{2} \times 8
    = \bigl(2.3\times10^{7}\bigr)^{2}\times 8
    \;\approx\; 4.232 \times 10^{15}\ \text{bytes},
\end{equation}
i.e.\ on the order of petabytes, which is plainly infeasible. In the
factorized form, by contrast, we store only $U$, $V$, and the core $W$:
\begin{equation}
  \underbrace{2\,(N \times r)\times 8}_{U,\,V}
   + \underbrace{r^{2}\times 8}_{W}
    \;=\; 2\,(2.3\times10^{7}\times 100)\times 8 + \text{(negligible)}
    \;\approx\; 36.8\ \text{GB}.
\end{equation}
The reduction is from petabytes to tens of gigabytes --- five orders of
magnitude --- and brings the $K$-step computation within reach of a single
large-memory machine. Table~\ref{tab:mf-memory} summarises the comparison.

\begin{table}[ht]
  \centering
  \caption{Memory footprint of the $K$-step transition computation
           ($N=2.3\times10^{7}$, $r=100$, double precision).}
  \label{tab:mf-memory}
  \begin{tabular}{lll}
    \toprule
    Representation & Quantity stored & Memory \\
    \midrule
    Dense $M^{k}$            & $N^{2}$ entries          & $\approx 4.232\times10^{15}$ B \\
    Factorized $U W^{k-1}V^{\top}$ & $2Nr + r^{2}$ entries & $\approx 36.8$ GB \\
    \bottomrule
  \end{tabular}
\end{table}

%--------------------------------------------------------------
\section{Factorization Methods}
\label{sec:mf-methods}

The reformulation of Section~\ref{sec:mf-solution} presumes a factorization
$M \approx U V^{\top}$. How that factorization is obtained is itself a design
choice, and three methods were considered: Alternating Least Squares (ALS),
Non-negative Matrix Factorization (NMF), and Weighted NMF. Because the entries
of $M$ are transition probabilities, they are non-negative by construction ---
a property that, as we shall see, decisively favours the non-negative methods.

\subsection{Alternating Least Squares (ALS)}
\label{sec:als}
ALS is a classical matrix-factorization technique, widely used in
collaborative filtering and recommender systems~\cite{HKV08,KBV09}. It
decomposes a large, sparse matrix into lower-dimensional latent factors by
minimising the regularised reconstruction error
\begin{equation}
  \label{eq:als-obj}
  \min_{U,\,V}\;
    \bigl\| M - U V^{\top} \bigr\|_{F}^{2}
    \;+\; \lambda\bigl(\|U\|_{F}^{2} + \|V\|_{F}^{2}\bigr),
\end{equation}
where $\lambda > 0$ is a regularisation weight and $\|\cdot\|_{F}$ the
Frobenius norm. The objective in \eqref{eq:als-obj} is not jointly convex in
$(U,V)$, but it \emph{is} convex in each factor with the other held fixed.
ALS exploits this by alternating: it fixes $V$ and solves a least-squares
problem for $U$, then fixes $U$ and solves for $V$, repeating until
convergence. Each subproblem has a closed form,
\begin{align}
  \label{eq:als-update}
  u_{i} &= \bigl(V^{\top}V + \lambda I_{r}\bigr)^{-1} V^{\top} m_{i},
   &\text{(row $i$ of $U$)} \nonumber\\
  v_{j} &= \bigl(U^{\top}U + \lambda I_{r}\bigr)^{-1} U^{\top} m_{j},
   &\text{(row $j$ of $V$)}
\end{align}
where $m_{i}$ and $m_{j}$ denote the relevant row/column of $M$ and $I_{r}$ is
the $r\times r$ identity. ALS is efficient, scalable, parallelises naturally
across rows, and copes well with sparse data, which is why it is a standard
first choice for recommendation at scale.

\paragraph{Why ALS does not work here.}
Despite its appeal, ALS proved unsuitable for this problem. The entries of $M$
are transition \emph{probabilities} and are therefore non-negative by
definition; a faithful factorization should respect this. ALS, however, places
no sign constraint on $U$ or $V$, and the reconstructed $UV^{\top}$ readily
produces negative entries that have no probabilistic meaning. In the
experiments conducted, approximately $34\%$ of the reconstructed entries were
negative, and the normalised reconstruction error was $51.23\%$. Both figures
are unacceptable for a probability matrix. This observation is what motivates
moving to non-negativity-constrained factorizations.

\subsection{Non-Negative Matrix Factorization (NMF)}
\label{sec:nmf}
NMF decomposes a non-negative matrix $M$ into two lower-dimensional factors
that are \emph{themselves} constrained to be non-negative,
$M \approx U V^{\top}$ with $U, V \ge 0$ elementwise. The objective is
\begin{equation}
  \label{eq:nmf-obj}
  \min_{U,\,V \,\ge\, 0}\;
    \bigl\| M - U V^{\top} \bigr\|_{F}^{2}.
\end{equation}
The non-negativity constraint guarantees that the reconstruction
$U V^{\top}$ is itself non-negative, so every entry remains interpretable as a
(non-normalised) probability --- precisely the property ALS failed to
preserve. The standard solver is the multiplicative-update rule of Lee and
Seung~\cite{LeeSeung01}, which preserves non-negativity automatically at every
iteration:
\begin{equation}
  \label{eq:nmf-update}
  U \;\leftarrow\; U \;\odot\; \frac{M V}{\,U V^{\top} V\,},
  \qquad
  V \;\leftarrow\; V \;\odot\; \frac{M^{\top} U}{\,V U^{\top} U\,},
\end{equation}
where $\odot$ denotes the Hadamard (elementwise) product and the division is
elementwise. Starting from non-negative initial factors, the updates in
\eqref{eq:nmf-update} keep $U$ and $V$ non-negative throughout and
monotonically decrease the objective \eqref{eq:nmf-obj}.

\subsection{Weighted NMF (WNMF)}
\label{sec:wnmf}
Weighted NMF extends NMF by attaching a weight to each entry of $M$, so that
different observations contribute with different importance. This is valuable
when the data contains missing values, noise, or entries of varying
confidence --- all of which characterise an extremely sparse association
matrix. Let $W^{\mathrm{wt}} \in \mathbb{R}^{N\times N}_{\ge 0}$ be the weight
matrix. The objective becomes
\begin{equation}
  \label{eq:wnmf-obj}
  \min_{U,\,V \,\ge\, 0}\;
    \bigl\| W^{\mathrm{wt}} \odot \bigl(M - U V^{\top}\bigr) \bigr\|_{F}^{2},
\end{equation}
and the corresponding multiplicative updates are
\begin{equation}
  \label{eq:wnmf-update}
  U \;\leftarrow\; U \;\odot\;
    \frac{\bigl(W^{\mathrm{wt}}\odot M\bigr) V}
         {\bigl(W^{\mathrm{wt}}\odot (U V^{\top})\bigr) V},
  \qquad
  V \;\leftarrow\; V \;\odot\;
    \frac{\bigl(W^{\mathrm{wt}}\odot M\bigr)^{\top} U}
         {\bigl(W^{\mathrm{wt}}\odot (U V^{\top})\bigr)^{\top} U}.
\end{equation}
Setting $W^{\mathrm{wt}} = \mathbf{1}$ (all weights equal) recovers ordinary
NMF; the freedom lies in choosing $W^{\mathrm{wt}}$ to reflect the structure of
the data.

%--------------------------------------------------------------
\section{Results and Discussion}
\label{sec:mf-results}

The decisive design lever turned out to be the choice of weights in WNMF. By
taking \emph{binary} weights --- $W^{\mathrm{wt}}_{ij} = 1$ where $M_{ij}$ is
observed and $0$ where it is missing --- the factorization is asked to fit only
the entries that actually carry information, and the vast field of structural
zeros in the sparse matrix is prevented from dominating the loss. This removes
the undesired influence of missing elements on the factors. Under binary
weighting the reconstruction error on the observed entries dropped almost to
zero, a dramatic improvement over both ALS (Section~\ref{sec:als}) and
unweighted NMF.

% FIGURE PLACEHOLDER: user to supply the error-comparison plot (ALS vs NMF
% vs binary-weighted WNMF) into figures/.
% \begin{figure}[ht]
%   \centering
%   \includegraphics[width=0.8\textwidth]{figures/mf_error_comparison.png}
%   \caption{Normalised reconstruction error across the three factorization
%            methods.}
%   \label{fig:mf-error}
% \end{figure}

Table~\ref{tab:mf-comparison} collects the comparison qualitatively; the
quantitative error figures from the experiments are reported where available.

\begin{table}[ht]
  \centering
  \caption{Comparison of the three factorization methods for the non-negative,
           highly sparse transition matrix.}
  \label{tab:mf-comparison}
  \begin{tabular}{lccc}
    \toprule
    Method & Non-negativity & Handles missing entries & Reconstruction error \\
    \midrule
    ALS              & No (\,$\approx34\%$ negative\,) & No   & $51.23\%$ \\
    NMF              & Yes                              & No   & --- \\
    WNMF (binary $W$) & Yes                             & Yes  & $\approx 0$ \\
    \bottomrule
  \end{tabular}
\end{table}

\noindent The progression across the three methods tells a clear story. ALS,
the conventional choice, breaks on the non-negativity requirement intrinsic to
a probability matrix. NMF restores non-negativity but still tries to fit the
structural zeros of a $10^{-4}$-dense matrix. Binary-weighted NMF addresses
both issues at once and drives the error on the observed entries almost to
zero, which is why it was adopted as the working method at the end of the
mid-term stage. The closer evaluation carried out subsequently, however, told a
less flattering story on two fronts. Empirically, recommendations generated from
the factorized matrix did not match the incumbent frequency baseline
(Section~\ref{sec:mf-ctr}); and structurally, the near-zero error figure was
found to conceal a serious failure mode (Section~\ref{sec:wnmf-limit}). The next
two sections take these in turn.

%--------------------------------------------------------------
\section{Recommendation-Quality Evaluation of the Factorization}
\label{sec:mf-ctr}

The reconstruction error of Section~\ref{sec:mf-results} measures how faithfully
the factors reproduce the \emph{observed} transitions. It says nothing directly
about the quality of the recommendations produced from the factorized matrix.
This section closes that gap: it evaluates the non-negative factorization ---
optimised over the observed transitions as in Section~\ref{sec:wnmf} --- on the
task it is actually meant for, ranking next products, and compares it against the
production frequency baseline. Three questions are asked in turn. Does the
factorized matrix recommend as well as the raw transition matrix and the
production model? Does the factorized representation preserve quality when it is
used for its intended purpose, cheap higher-order transitions? And does adding a
self-transition term to the matrix help the factorization close the gap?

\subsection{Evaluation Setup}
\label{sec:mf-ctr-setup}

\paragraph{Subcategory-wise factorization.}
Factorising the complete transition matrix was not feasible within the available
hardware. At the scale of Section~\ref{sec:mf-problem}, storing the factors and
optimising the reconstruction objective across the full catalogue --- and doing
so repeatedly while sweeping factorization ranks --- exceeded the memory budget.
The evaluation was therefore carried out one product subcategory at a time. For
each selected subcategory a row-normalised transition probability matrix was
built using only the products belonging to that subcategory. These per-%
subcategory matrices are far smaller than the global matrix while preserving the
transition structure within their product domain, which makes it possible to
assess both reconstruction quality and its effect on recommendations without the
memory obstacle.

\paragraph{Recommendation generation.}
Recommendations were read off the reconstructed transition matrix by taking, for
each source product, the products with the highest transition probability in its
row. The top~$20$ such products were kept as the candidate recommendations for
that source.

\paragraph{Backfilling.}
In practice not every product had a full list of twenty recommendations once its
transition probabilities were ranked. To keep every list at exactly twenty
items, any empty slots were filled with the recommendations already produced by
IndiaMART's existing production system --- a \emph{backfilling} step. This keeps
the list length uniform across methods and keeps the output compatible with the
production pipeline. Backfilling recurs as a baseline component in later chapters
(see Section~\ref{sec:emb-baselines}).

\paragraph{Baseline.}
The comparison point is the \emph{Frequency-Based Model} (FBM), the production
recommender. FBM ranks candidates directly by the transition probabilities of
the (row-normalised) transition matrix, with no factorization: the products with
the highest observed transition probability are recommended first. The specific
production version appearing in the tables below is denoted FBM~v1.5.1.

\paragraph{Metric.}
Recommendation quality is measured by click-through rate (CTR), the fraction of
recommended products that users clicked after being shown them. CTR is reported
separately for the IndiaMART website (IM) and the IndiaMART mobile platform
(IMOB).

\subsection{Factorization versus the Transition Matrix and FBM}
\label{sec:mf-ctr-subcat}

Table~\ref{tab:mf-ctr-subcat} compares three recommenders across six
subcategories: recommendations from the NMF-reconstructed matrix, recommendations
from the raw transition matrix with backfilling, and the production FBM.

\begin{table}[ht]
  \centering
  \caption{CTR (\%) by subcategory: NMF reconstruction versus the raw transition
           matrix with backfilling versus the production FBM~v1.5.1.}
  \label{tab:mf-ctr-subcat}
  \small
  \begin{tabular}{lcccccc}
    \toprule
    & \multicolumn{2}{c}{NMF} & \multicolumn{2}{c}{Transition + backfilling}
    & \multicolumn{2}{c}{FBM v1.5.1} \\
    \cmidrule(lr){2-3}\cmidrule(lr){4-5}\cmidrule(lr){6-7}
    Subcat & IM & IMOB & IM & IMOB & IM & IMOB \\
    \midrule
    9   & $8.52$ & $11.73$ & $14.30$ & $16.33$ & $14.41$ & $15.31$ \\
    22  & $9.55$ & $15.52$ & $12.73$ & $19.34$ & $14.28$ & $19.95$ \\
    144 & $5.28$ & $7.46$  & $12.41$ & $15.32$ & $9.46$  & $11.44$ \\
    407 & $7.62$ & $8.78$  & $11.61$ & $13.72$ & $10.05$ & $13.89$ \\
    487 & $5.32$ & $9.83$  & $11.11$ & $17.97$ & $10.18$ & $18.35$ \\
    799 & $6.53$ & $10.57$ & $12.72$ & $15.50$ & $11.59$ & $15.80$ \\
    \bottomrule
  \end{tabular}
\end{table}

The pattern is the same in every subcategory. The NMF reconstruction produces the
lowest CTR of the three, on both surfaces, without exception. Recommendations
taken directly from the transition matrix with backfilling, by contrast, are
roughly on par with the production system --- comparable to, and in several
subcategories slightly above, FBM~v1.5.1. The exact transition probabilities
therefore carry more useful signal than their low-rank approximation does. A
plausible reason is that the transition matrix is extremely sparse and its useful
structure is highly localised: a rank-$r$ non-negative factorization compresses
the matrix aggressively, and the approximation smooths away precisely the sharp,
product-specific transitions that a good next-product recommendation depends on.

\subsection{Higher-Order Transitions With and Without Factorization}
\label{sec:mf-ctr-higher}

The comparison above evaluates the factorization only through its first-order
reconstruction. But reconstructing the matrix was never the point; the reason for
factorising at all (Section~\ref{sec:mf-solution}) was to compute higher-order
transitions cheaply, using the factored-power identity~\eqref{eq:factored-power}
in which the expensive part reduces to powering the small $r\times r$ core. The
question is whether recommendations built from those factorized higher-order
transitions hold up against the same transitions computed by directly powering
the raw matrix.

For each subcategory two pipelines were run for exponents $k\in\{1,2,3,4,5\}$: in
the first the raw transition matrix was repeatedly multiplied to form $M^{k}$; in
the second the matrix was factorised and the higher-order transitions were
obtained through the low-rank representation. In both cases recommendations were
generated by ranking each row of the resulting matrix, and CTR was measured on
both surfaces. Table~\ref{tab:mf-ctr-higher} collects the results.

\begin{table}[ht]
  \centering
  \caption{Higher-order transition CTR (\%) without and with factorization, by
           subcategory and transition exponent $k$.}
  \label{tab:mf-ctr-higher}
  \small
  \begin{tabular}{clcccc}
    \toprule
    & & \multicolumn{2}{c}{Without factorization}
      & \multicolumn{2}{c}{With factorization} \\
    \cmidrule(lr){3-4}\cmidrule(lr){5-6}
    Subcat & $k$ & IM & IMOB & IM & IMOB \\
    \midrule
    \multirow{5}{*}{9}
      & 1 & $7.11$ & $10.63$ & $8.52$ & $11.13$ \\
      & 2 & $6.52$ & $9.15$  & $6.07$ & $9.46$ \\
      & 3 & $7.26$ & $12.02$ & $8.30$ & $12.23$ \\
      & 4 & $6.96$ & $10.46$ & $6.96$ & $10.27$ \\
      & 5 & $7.85$ & $11.40$ & $7.41$ & $10.89$ \\
    \midrule
    \multirow{5}{*}{22}
      & 1 & $8.89$  & $15.25$ & $9.55$  & $14.99$ \\
      & 2 & $9.21$  & $14.40$ & $8.25$  & $13.85$ \\
      & 3 & $10.79$ & $16.60$ & $10.16$ & $14.45$ \\
      & 4 & $8.57$  & $13.99$ & $8.57$  & $11.48$ \\
      & 5 & $8.57$  & $13.00$ & $6.67$  & $10.19$ \\
    \midrule
    \multirow{5}{*}{144}
      & 1 & $4.19$ & $6.90$ & $5.28$ & $7.23$ \\
      & 2 & $4.42$ & $6.06$ & $3.95$ & $6.38$ \\
      & 3 & $6.05$ & $8.20$ & $6.05$ & $8.24$ \\
      & 4 & $6.28$ & $6.18$ & $6.28$ & $6.66$ \\
      & 5 & $5.58$ & $7.83$ & $5.35$ & $7.27$ \\
    \midrule
    \multirow{5}{*}{407}
      & 1 & $5.38$ & $8.02$ & $7.62$ & $8.24$ \\
      & 2 & $6.05$ & $6.67$ & $5.55$ & $6.98$ \\
      & 3 & $7.35$ & $9.56$ & $6.39$ & $8.83$ \\
      & 4 & $5.88$ & $8.02$ & $5.71$ & $7.47$ \\
      & 5 & $7.06$ & $9.34$ & $5.55$ & $8.22$ \\
    \midrule
    \multirow{5}{*}{487}
      & 1 & $4.29$ & $7.16$ & $5.32$ & $9.26$ \\
      & 2 & $3.06$ & $6.82$ & $3.98$ & $7.35$ \\
      & 3 & $5.97$ & $8.97$ & $6.74$ & $9.02$ \\
      & 4 & $4.44$ & $8.35$ & $5.05$ & $8.44$ \\
      & 5 & $6.89$ & $9.06$ & $6.28$ & $9.02$ \\
    \midrule
    \multirow{5}{*}{799}
      & 1 & $5.81$ & $10.81$ & $6.53$ & $11.42$ \\
      & 2 & $6.59$ & $10.29$ & $6.82$ & $10.19$ \\
      & 3 & $7.82$ & $12.33$ & $8.16$ & $11.78$ \\
      & 4 & $7.15$ & $11.62$ & $7.49$ & $11.10$ \\
      & 5 & $7.82$ & $11.65$ & $7.60$ & $10.58$ \\
    \bottomrule
  \end{tabular}
\end{table}

Two things emerge. At the first order ($k=1$) the factorized representation is
competitive, and in several subcategories (9, 144, 407, 487, 799) it is even a
little ahead of the raw matrix --- confirming that the approximation captures the
dominant direct transitions with only small error. As the exponent grows,
however, the factorized results become erratic. In subcategories such as 22 and
407 the raw matrix pulls clearly ahead at larger $k$, while in others the two stay
close; no subcategory shows the factorization improving with order. The most
natural reading is that each multiplication of the approximate matrix compounds
its reconstruction error, so the higher-order transition probabilities --- and
hence the rankings drawn from them --- drift further from the true chain as $k$
increases. Whether factorization helps or hurts at higher order also varies with
the subcategory, which suggests the effect depends on the structure of the
underlying transition graph rather than being a property of the method alone.

Crucially, the earlier comparison still stands over the top of all this. Even
where the factorized first-order result edges out the raw matrix, it remains below
the production FBM~v1.5.1 of Table~\ref{tab:mf-ctr-subcat}. An improvement over
the raw transition matrix is of little value if it does not also clear the
production baseline, and here it does not.

\subsection{Incorporating a Self-Transition Probability}
\label{sec:mf-ctr-self}

A final experiment revisits an assumption built into the transition matrix. In
all of the above, the diagonal of the matrix was zero: a product could never
transition to itself. This is convenient but ignores a real browsing behaviour
--- users often linger on a product, revisiting its details, comparing
specifications, or returning to it after looking at related items, before moving
on. A non-zero diagonal lets the Markov model represent this.

Concretely, let $M=[m_{ij}]$ be the original row-normalised matrix with
$m_{ii}=0$ and $\sum_j m_{ij}=1$. A fixed self-transition probability $\alpha$ is
placed on the diagonal and the off-diagonal mass is rescaled to keep each row
stochastic:
\begin{equation}
  \label{eq:self-transition}
  m'_{ii} = \alpha,
  \qquad
  m'_{ij} = (1-\alpha)\,m_{ij}\ \ (i\neq j),
  \qquad\text{so that}\quad \sum_{j} m'_{ij} = 1 .
\end{equation}
The experiment then sweeps three factors jointly --- the self-transition
probability $\alpha$, the factorization rank, and the transition exponent --- on
subcategory~526. Factorization ranks of $1000$, $1500$, and $2000$ are compared
against recommendations taken directly from the (self-transition-augmented) raw
matrix. Table~\ref{tab:mf-ctr-self} reports the results.

\begin{table}[ht]
  \centering
  \caption{CTR (\%) on subcategory~526 with a self-transition probability
           $\alpha$, for factorization ranks $1000/1500/2000$ and for the raw
           matrix (``No fac.''), across transition exponents $k$.}
  \label{tab:mf-ctr-self}
  \small
  \resizebox{\textwidth}{!}{%
  \begin{tabular}{clcccccccc}
    \toprule
    & & \multicolumn{2}{c}{Rank 1000} & \multicolumn{2}{c}{Rank 1500}
      & \multicolumn{2}{c}{Rank 2000} & \multicolumn{2}{c}{No fac.} \\
    \cmidrule(lr){3-4}\cmidrule(lr){5-6}\cmidrule(lr){7-8}\cmidrule(lr){9-10}
    $\alpha$ & $k$ & IM & IMOB & IM & IMOB & IM & IMOB & IM & IMOB \\
    \midrule
    \multirow{4}{*}{$0.0625$}
      & 1 & $13.51$ & $22.67$ & $15.32$ & $27.12$ & $15.74$ & $28.26$ & $15.32$ & $28.32$ \\
      & 2 & $13.23$ & $20.97$ & $14.76$ & $25.49$ & $16.43$ & $26.30$ & $15.60$ & $26.52$ \\
      & 3 & $12.26$ & $18.09$ & $13.09$ & $21.56$ & $13.93$ & $22.45$ & $14.21$ & $23.27$ \\
      & 4 & $10.58$ & $16.31$ & $11.00$ & $17.61$ & $11.56$ & $18.90$ & $11.56$ & $19.69$ \\
    \midrule
    \multirow{4}{*}{$0.125$}
      & 1 & $14.21$ & $24.25$ & $15.18$ & $27.22$ & $16.57$ & $27.95$ & $15.32$ & $28.32$ \\
      & 2 & $14.90$ & $22.83$ & $15.88$ & $26.75$ & $16.57$ & $27.53$ & $16.30$ & $27.56$ \\
      & 3 & $12.67$ & $20.06$ & $14.35$ & $23.57$ & $14.62$ & $23.90$ & $14.62$ & $24.44$ \\
      & 4 & $11.28$ & $17.01$ & $11.56$ & $18.75$ & $12.12$ & $19.64$ & $12.40$ & $21.05$ \\
    \midrule
    \multirow{4}{*}{$0.25$}
      & 1 & $14.76$ & $24.67$ & $15.46$ & $26.58$ & $16.57$ & $27.26$ & $15.32$ & $28.32$ \\
      & 2 & $15.60$ & $24.33$ & $15.18$ & $26.68$ & $16.85$ & $27.53$ & $16.71$ & $28.44$ \\
      & 3 & $13.93$ & $22.52$ & $14.35$ & $24.93$ & $15.32$ & $25.80$ & $16.43$ & $26.71$ \\
      & 4 & $12.26$ & $19.36$ & $13.23$ & $21.47$ & $13.65$ & $22.37$ & $14.48$ & $24.14$ \\
    \midrule
    \multirow{4}{*}{$0.5$}
      & 1 & $14.76$ & $23.97$ & $14.76$ & $24.54$ & $14.76$ & $25.57$ & $15.32$ & $28.32$ \\
      & 2 & $15.60$ & $24.89$ & $15.18$ & $25.89$ & $15.74$ & $26.66$ & $16.85$ & $28.86$ \\
      & 3 & $14.90$ & $24.65$ & $14.90$ & $26.02$ & $16.02$ & $26.40$ & $17.55$ & $28.41$ \\
      & 4 & $13.79$ & $23.44$ & $14.21$ & $25.16$ & $15.46$ & $25.35$ & $16.99$ & $27.63$ \\
    \bottomrule
  \end{tabular}%
  }
\end{table}

The results are internally sensible but do not change the conclusion. Within the
factorized models, a higher rank almost always helps: rank~2000 is consistently
the strongest of the three, as expected, since a larger rank retains more of the
transition information. Recommendation quality also falls off steadily as the
exponent grows, for both the factorized and the raw matrix and at every value of
$\alpha$ --- higher-order powers spread probability mass over ever more products
and dilute the strong direct transitions that drive good recommendations. The
value of $\alpha$ trades off how much probability is held at the current product
against how much flows to its neighbours, and moderate values keep that balance
best. But across the whole sweep --- every rank, every exponent, every $\alpha$
--- the raw self-transition matrix remains competitive and, in most cells,
achieves the highest CTR overall, especially at the larger exponents. Adding a
self-transition term does not let the factorized representation overcome the
information it loses to the low-rank approximation.

\subsection{Summary}
\label{sec:mf-ctr-summary}

The three experiments agree. Non-negative matrix factorization gives a compact,
computationally cheap way to approximate the transition matrix and its powers,
but it does not preserve recommendation quality: its first-order recommendations
trail the raw matrix, its higher-order recommendations accumulate approximation
error, and a self-transition term does not close the gap. In every setting the
factorized recommender stays below the production FBM. Because beating the raw
transition matrix is not enough --- the method has to beat the production
baseline to be worth deploying --- the factorization approach was not carried
forward as the final recommender. This empirical shortfall, together with the
structural limitation examined next, is what motivated the alternative directions
pursued in the following chapters.

%--------------------------------------------------------------
\section{Limitation of Weighted NMF}
\label{sec:wnmf-limit}

The very mechanism that makes binary-weighted NMF attractive --- fitting only
the observed entries --- is also the source of its weakness. Recall from
\eqref{eq:wnmf-obj} that with a binary weight matrix the loss is accumulated
\emph{only} over positions where $W^{\mathrm{wt}}_{ij}=1$. The multiplicative
updates of \eqref{eq:wnmf-update} therefore move $U$ and $V$ so as to improve
the reconstruction of those observed positions and nothing else. Likewise, the
absolute (and normalised) error is measured only against the observed entries
of the original matrix. The reconstructed values at the \emph{unobserved}
positions --- those with $W^{\mathrm{wt}}_{ij}=0$ --- enter neither the loss nor
the error metric. They are produced as an uncontrolled by-product of the
factors and are never monitored during training.

This is precisely where the method fails. Because nothing in the objective
constrains the unobserved entries, there is no force keeping them small or
well-behaved. As the factors are driven to fit the observed entries, some of
the unconstrained entries in the reconstruction $U V^{\top}$ drift to
unusually large values. Two consequences follow:

\begin{itemize}
  \item \textbf{The error metric is blind to the problem.} Since error is
        computed only on observed entries, a model can report near-zero error
        while simultaneously placing wildly inflated values at unobserved
        positions. The optimistic figure of Section~\ref{sec:mf-results} is, in
        this light, partly an artefact of \emph{what} is being measured.
  \item \textbf{Inflated entries dominate the Markov chain.} In the $K$-step
        recommender these spuriously large reconstructed values behave as if
        they were high transition probabilities. Propagated through
        $M^{K}\approx U\,W^{\,K-1}\,V^{\top}$, they dominate the chain's
        transitions, so that a handful of products surface excessively in the
        recommendations regardless of the query --- crowding out the genuinely
        relevant items.
\end{itemize}

In short, binary-weighted NMF optimises a quantity (observed-entry
reconstruction) that is not aligned with the quantity that actually matters for
recommendation (the behaviour of the \emph{full} reconstructed matrix under
repeated multiplication). Controlling the unobserved entries --- so that the
reconstruction remains a well-behaved transition operator everywhere, not only
where data was seen --- is the open problem that this work carries into the
next stage.
